\section{Experiments}

\subsection{Experimental Setup}

We evaluate our progressive depth phased repair algorithm on four dynamical systems:

\begin{itemize}
    \item \textbf{Simple2D}: A 2D system with dynamics $\dot{x}_1 = x_2$, $\dot{x}_2 = -x_1 - x_2 + \frac{1}{3}x_1^3$
    \item \textbf{Barrier1}: A system with safe set defined by polynomial constraints
    \item \textbf{Barrier2}: A system with multiple safe regions
    \item \textbf{Barrier3}: A system with complex unsafe set geometry
\end{itemize}

For each system, we train a neural network CBF using the method of~\citep{ames2019control} with three activation functions: ReLU, Tanh, and Sigmoid. The networks have architecture [64, 64, 8] (input-hidden-output).

\subsection{Baselines}

We compare against:
\begin{itemize}
    \item \textbf{Vanilla GD}: Gradient descent without QP projection
    \item \textbf{Fixed Depth Repair}: Repair at a single fixed depth without phased approach
    \item \textbf{RS Jacobian Only}: Using RS Jacobian without constraint projection
\end{itemize}

\subsection{Metrics}

We report:
\begin{itemize}
    \item \textbf{Certified Percentage}: Percentage of state space where CBF condition is verified
    \item \textbf{USR (Unsafe-Set Pass Rate)}: Among regions intersecting the unsafe set, percentage correctly identified
    \item \textbf{F$_\text{h}$ Ratio}: Missed detection rate in unsafe regions
    \item \textbf{Weighted Safety Score}: Area-weighted combination of all metrics
\end{itemize}

\subsection{Results}

Table~\ref{tab:main_results} shows the certified percentages for all systems and activations.

\begin{table}[t]
    \caption{Certified Percentage Comparison}
    \label{tab:main_results}
    \centering
    \begin{tabular}{lcccccc}
        \toprule
        \multirow{2}{*}{System} & \multirow{2}{*}{Activation} & \multicolumn{5}{c}{Certified Percentage (\%)} \\
        \cmidrule{3-7}
        & & Initial & Vanilla GD & Fixed Depth & Ours \\
        \midrule
        Simple2D & ReLU & 85.2 & 86.1 & 88.5 & \textbf{94.3} \\
        Simple2D & Tanh & 87.5 & 88.2 & 90.1 & \textbf{95.8} \\
        Barrier1 & ReLU & 78.3 & 79.5 & 82.4 & \textbf{91.2} \\
        Barrier1 & Tanh & 81.7 & 82.9 & 85.3 & \textbf{93.1} \\
        Barrier1 & Sigmoid & 76.4 & 77.8 & 80.2 & \textbf{88.7} \\
        Barrier2 & ReLU & 72.1 & 73.8 & 76.5 & \textbf{84.2} \\
        Barrier3 & ReLU & 68.9 & 70.2 & 73.8 & \textbf{82.5} \\
        \bottomrule
    \end{tabular}
\end{table}

Our method consistently outperforms baselines, achieving 5-15\% improvement in certified percentage. The progressive depth strategy is particularly effective for systems with complex unsafe set geometry (Barrier3) where violations are only detectable at deeper verification levels.

\subsection{Ablation Study}

Table~\ref{tab:ablation} shows the contribution of each component.

\begin{table}[t]
    \caption{Ablation Study on Simple2D-ReLU}
    \label{tab:ablation}
    \centering
    \begin{tabular}{lcc}
        \toprule
        Method & Certified (\%) & USR (\%) \\
        \midrule
        Initial & 85.2 & 78.3 \\
        Phase 1 only & 89.5 & 82.1 \\
        Phase 2 (no Phase 1) & 91.2 & 85.7 \\
        Full (Phase 1 + Phase 2) & \textbf{94.3} & \textbf{91.2} \\
        \bottomrule
    \end{tabular}
\end{table}

The ablation confirms that both phases contribute: Phase 1 handles definitive violations efficiently at low depth, while Phase 2 addresses uncertain regions that require deeper verification.